\QHeading{第1题　画出把下面的平行四边形先向右平移7格，再向下平移5格后}

\par\noindent\makebox[\linewidth][c]{
\begin{tikzpicture}[x=1bp,y=-1bp,color=black,draw=black,text=black]
\begin{scope}\clip (56.0,115.0) rectangle (458.0,339.33);
\hdhead{78}{图形的运动——信息窗2}
\hdsection{1}{137}{画出把下面的平行四边形先向右平移7格，再向下平移5格后}
\q{117}{179}{699}{13.5}{\color{bookred}的图形。}
\hdsection{2}{512}{按要求画一画。}
\q{92}{558}{722}{12.5}{① 将图A向右平移14格后，再向下平移5格。}
\q{92}{601}{722}{12.5}{② 将图B向左平移9格后，再向上平移5格。}
\q{92}{645}{722}{12.5}{③ 将图C绕$O$点顺时针旋转$90^\circ$。}
\q{92}{688}{722}{12.5}{④ 将图D绕$O\prime$点逆时针旋转$90^\circ$。}
\begin{scope}[scale=.6666667]
\draw[hdgrid] (215,200)--(215,472.25);
\draw[hdgrid] (239.45,200)--(239.45,472.25);
\draw[hdgrid] (263.9,200)--(263.9,472.25);
\draw[hdgrid] (288.35,200)--(288.35,472.25);
\draw[hdgrid] (312.8,200)--(312.8,472.25);
\draw[hdgrid] (337.25,200)--(337.25,472.25);
\draw[hdgrid] (361.7,200)--(361.7,472.25);
\draw[hdgrid] (386.15,200)--(386.15,472.25);
\draw[hdgrid] (410.6,200)--(410.6,472.25);
\draw[hdgrid] (435.05,200)--(435.05,472.25);
\draw[hdgrid] (459.5,200)--(459.5,472.25);
\draw[hdgrid] (483.95,200)--(483.95,472.25);
\draw[hdgrid] (508.4,200)--(508.4,472.25);
\draw[hdgrid] (532.85,200)--(532.85,472.25);
\draw[hdgrid] (557.3,200)--(557.3,472.25);
\draw[hdgrid] (581.75,200)--(581.75,472.25);
\draw[hdgrid] (606.2,200)--(606.2,472.25);
\draw[hdgrid] (630.65,200)--(630.65,472.25);
\draw[hdgrid] (655.1,200)--(655.1,472.25);
\draw[hdgrid] (679.55,200)--(679.55,472.25);
\draw[hdgrid] (215,200)--(679.55,200);
\draw[hdgrid] (215,224.75)--(679.55,224.75);
\draw[hdgrid] (215,249.5)--(679.55,249.5);
\draw[hdgrid] (215,274.25)--(679.55,274.25);
\draw[hdgrid] (215,299)--(679.55,299);
\draw[hdgrid] (215,323.75)--(679.55,323.75);
\draw[hdgrid] (215,348.5)--(679.55,348.5);
\draw[hdgrid] (215,373.25)--(679.55,373.25);
\draw[hdgrid] (215,398)--(679.55,398);
\draw[hdgrid] (215,422.75)--(679.55,422.75);
\draw[hdgrid] (215,447.5)--(679.55,447.5);
\draw[hdgrid] (215,472.25)--(679.55,472.25);
\draw[hdoutline] (239.45,274.25)--(263.9,224.75)--(361.7,224.75)--(337.25,274.25)--cycle;
\draw[hdgrid] (137,733)--(137,1126.6);
\draw[hdgrid] (161.15,733)--(161.15,1126.6);
\draw[hdgrid] (185.3,733)--(185.3,1126.6);
\draw[hdgrid] (209.45,733)--(209.45,1126.6);
\draw[hdgrid] (233.6,733)--(233.6,1126.6);
\draw[hdgrid] (257.75,733)--(257.75,1126.6);
\draw[hdgrid] (281.9,733)--(281.9,1126.6);
\draw[hdgrid] (306.05,733)--(306.05,1126.6);
\draw[hdgrid] (330.2,733)--(330.2,1126.6);
\draw[hdgrid] (354.35,733)--(354.35,1126.6);
\draw[hdgrid] (378.5,733)--(378.5,1126.6);
\draw[hdgrid] (402.65,733)--(402.65,1126.6);
\draw[hdgrid] (426.8,733)--(426.8,1126.6);
\draw[hdgrid] (450.95,733)--(450.95,1126.6);
\draw[hdgrid] (475.1,733)--(475.1,1126.6);
\draw[hdgrid] (499.25,733)--(499.25,1126.6);
\draw[hdgrid] (523.4,733)--(523.4,1126.6);
\draw[hdgrid] (547.55,733)--(547.55,1126.6);
\draw[hdgrid] (571.7,733)--(571.7,1126.6);
\draw[hdgrid] (595.85,733)--(595.85,1126.6);
\draw[hdgrid] (620,733)--(620,1126.6);
\draw[hdgrid] (644.15,733)--(644.15,1126.6);
\draw[hdgrid] (668.3,733)--(668.3,1126.6);
\draw[hdgrid] (692.45,733)--(692.45,1126.6);
\draw[hdgrid] (716.6,733)--(716.6,1126.6);
\draw[hdgrid] (740.75,733)--(740.75,1126.6);
\draw[hdgrid] (764.9,733)--(764.9,1126.6);
\draw[hdgrid] (137,733)--(764.9,733);
\draw[hdgrid] (137,757.6)--(764.9,757.6);
\draw[hdgrid] (137,782.2)--(764.9,782.2);
\draw[hdgrid] (137,806.8)--(764.9,806.8);
\draw[hdgrid] (137,831.4)--(764.9,831.4);
\draw[hdgrid] (137,856)--(764.9,856);
\draw[hdgrid] (137,880.6)--(764.9,880.6);
\draw[hdgrid] (137,905.2)--(764.9,905.2);
\draw[hdgrid] (137,929.8)--(764.9,929.8);
\draw[hdgrid] (137,954.4)--(764.9,954.4);
\draw[hdgrid] (137,979)--(764.9,979);
\draw[hdgrid] (137,1003.6)--(764.9,1003.6);
\draw[hdgrid] (137,1028.2)--(764.9,1028.2);
\draw[hdgrid] (137,1052.8)--(764.9,1052.8);
\draw[hdgrid] (137,1077.4)--(764.9,1077.4);
\draw[hdgrid] (137,1102)--(764.9,1102);
\draw[hdgrid] (137,1126.6)--(764.9,1126.6);
\draw[hdoutline] (281.9,757.6)--(330.2,806.8)--(281.9,856)--(233.6,806.8)--cycle;
\node[font=\normalsize,inner sep=0pt] at (293.975,816.64) {A};
\draw[hdoutline] (547.55,1028.2)--(595.85,1077.4)--(644.15,1028.2)--(644.15,1102)--(547.55,1102)--cycle;
\node[font=\normalsize,inner sep=0pt] at (607.925,1089.7) {B};
\draw[hdoutline] (185.3,954.4)--(185.3,905.2)--(281.9,905.2)--(233.6,954.4)--cycle;
\node[font=\normalsize,inner sep=0pt] at (223.94,942.1) {C};
\draw[hdoutline] (450.95,806.8)--(475.1,831.4)--(475.1,880.6)--(450.95,905.2)--(426.8,880.6)--(426.8,831.4)--cycle;
\node[font=\normalsize,inner sep=0pt] at (463.025,868.3) {D};
\fill[black] (185.3,954.4) circle (1.7);
\node[font=\normalsize,inner sep=0pt] at (171,970) {$O$};
\fill[black] (450.95,806.8) circle (1.7);
\node[font=\normalsize,inner sep=0pt] at (460,794) {$O\prime$};
\end{scope}

\begin{scope}[color=answerred,draw=answerred,text=answerred]\draw[answerred,line width=1pt] (274.07,265.40) -- (290.37,232.47) -- (355.57,232.47) -- (339.27,265.40) -- cycle;
\draw[answerred,line width=1pt] (412.67,587.33) -- (444.67,620.00) -- (412.67,653.33) -- (380.00,620.00) -- cycle;
\draw[answerred,line width=1pt] (219.33,604.00) -- (251.33,636.67) -- (283.33,604.00) -- (283.33,653.33) -- (219.33,653.33) -- cycle;
\draw[answerred,line width=1pt] (122.67,636.67) -- (154.67,636.67) -- (154.67,702.00) -- (122.67,669.33) -- cycle;
\draw[answerred,line width=1pt] (299.33,538.00) -- (315.33,522.00) -- (347.33,522.00) -- (363.33,538.00) -- (347.33,554.67) -- (315.33,554.67) -- cycle;
\ans{405.33}{612.67}{26.00}{11}{$A\prime$}
\ans{242.00}{636.67}{26.00}{11}{$B\prime$}
\ans{133.33}{670.67}{26.00}{11}{$C\prime$}
\ans{328.67}{531.33}{26.00}{11}{$D\prime$}
\end{scope}
\end{scope}\pgfresetboundingbox\path[use as bounding box] (56.0,115.0) rectangle (458.0,339.33);
\end{tikzpicture}}\par
\par\noindent




\par

\par\noindent

\QHeading{第2题　按要求画一画。}

\par\noindent\makebox[\linewidth][c]{
\begin{tikzpicture}[x=1bp,y=-1bp,color=black,draw=black,text=black]
\begin{scope}\clip (86.0,483.0) rectangle (515.0,756.0);
\hdhead{78}{图形的运动——信息窗2}
\hdsection{1}{137}{画出把下面的平行四边形先向右平移7格，再向下平移5格后}
\q{117}{179}{699}{13.5}{\color{bookred}的图形。}
\hdsection{2}{512}{按要求画一画。}




\begin{scope}[scale=.6666667]
\draw[hdgrid] (215,200)--(215,472.25);
\draw[hdgrid] (239.45,200)--(239.45,472.25);
\draw[hdgrid] (263.9,200)--(263.9,472.25);
\draw[hdgrid] (288.35,200)--(288.35,472.25);
\draw[hdgrid] (312.8,200)--(312.8,472.25);
\draw[hdgrid] (337.25,200)--(337.25,472.25);
\draw[hdgrid] (361.7,200)--(361.7,472.25);
\draw[hdgrid] (386.15,200)--(386.15,472.25);
\draw[hdgrid] (410.6,200)--(410.6,472.25);
\draw[hdgrid] (435.05,200)--(435.05,472.25);
\draw[hdgrid] (459.5,200)--(459.5,472.25);
\draw[hdgrid] (483.95,200)--(483.95,472.25);
\draw[hdgrid] (508.4,200)--(508.4,472.25);
\draw[hdgrid] (532.85,200)--(532.85,472.25);
\draw[hdgrid] (557.3,200)--(557.3,472.25);
\draw[hdgrid] (581.75,200)--(581.75,472.25);
\draw[hdgrid] (606.2,200)--(606.2,472.25);
\draw[hdgrid] (630.65,200)--(630.65,472.25);
\draw[hdgrid] (655.1,200)--(655.1,472.25);
\draw[hdgrid] (679.55,200)--(679.55,472.25);
\draw[hdgrid] (215,200)--(679.55,200);
\draw[hdgrid] (215,224.75)--(679.55,224.75);
\draw[hdgrid] (215,249.5)--(679.55,249.5);
\draw[hdgrid] (215,274.25)--(679.55,274.25);
\draw[hdgrid] (215,299)--(679.55,299);
\draw[hdgrid] (215,323.75)--(679.55,323.75);
\draw[hdgrid] (215,348.5)--(679.55,348.5);
\draw[hdgrid] (215,373.25)--(679.55,373.25);
\draw[hdgrid] (215,398)--(679.55,398);
\draw[hdgrid] (215,422.75)--(679.55,422.75);
\draw[hdgrid] (215,447.5)--(679.55,447.5);
\draw[hdgrid] (215,472.25)--(679.55,472.25);
\draw[hdoutline] (239.45,274.25)--(263.9,224.75)--(361.7,224.75)--(337.25,274.25)--cycle;
\draw[hdgrid] (137,733)--(137,1126.6);
\draw[hdgrid] (161.15,733)--(161.15,1126.6);
\draw[hdgrid] (185.3,733)--(185.3,1126.6);
\draw[hdgrid] (209.45,733)--(209.45,1126.6);
\draw[hdgrid] (233.6,733)--(233.6,1126.6);
\draw[hdgrid] (257.75,733)--(257.75,1126.6);
\draw[hdgrid] (281.9,733)--(281.9,1126.6);
\draw[hdgrid] (306.05,733)--(306.05,1126.6);
\draw[hdgrid] (330.2,733)--(330.2,1126.6);
\draw[hdgrid] (354.35,733)--(354.35,1126.6);
\draw[hdgrid] (378.5,733)--(378.5,1126.6);
\draw[hdgrid] (402.65,733)--(402.65,1126.6);
\draw[hdgrid] (426.8,733)--(426.8,1126.6);
\draw[hdgrid] (450.95,733)--(450.95,1126.6);
\draw[hdgrid] (475.1,733)--(475.1,1126.6);
\draw[hdgrid] (499.25,733)--(499.25,1126.6);
\draw[hdgrid] (523.4,733)--(523.4,1126.6);
\draw[hdgrid] (547.55,733)--(547.55,1126.6);
\draw[hdgrid] (571.7,733)--(571.7,1126.6);
\draw[hdgrid] (595.85,733)--(595.85,1126.6);
\draw[hdgrid] (620,733)--(620,1126.6);
\draw[hdgrid] (644.15,733)--(644.15,1126.6);
\draw[hdgrid] (668.3,733)--(668.3,1126.6);
\draw[hdgrid] (692.45,733)--(692.45,1126.6);
\draw[hdgrid] (716.6,733)--(716.6,1126.6);
\draw[hdgrid] (740.75,733)--(740.75,1126.6);
\draw[hdgrid] (764.9,733)--(764.9,1126.6);
\draw[hdgrid] (137,733)--(764.9,733);
\draw[hdgrid] (137,757.6)--(764.9,757.6);
\draw[hdgrid] (137,782.2)--(764.9,782.2);
\draw[hdgrid] (137,806.8)--(764.9,806.8);
\draw[hdgrid] (137,831.4)--(764.9,831.4);
\draw[hdgrid] (137,856)--(764.9,856);
\draw[hdgrid] (137,880.6)--(764.9,880.6);
\draw[hdgrid] (137,905.2)--(764.9,905.2);
\draw[hdgrid] (137,929.8)--(764.9,929.8);
\draw[hdgrid] (137,954.4)--(764.9,954.4);
\draw[hdgrid] (137,979)--(764.9,979);
\draw[hdgrid] (137,1003.6)--(764.9,1003.6);
\draw[hdgrid] (137,1028.2)--(764.9,1028.2);
\draw[hdgrid] (137,1052.8)--(764.9,1052.8);
\draw[hdgrid] (137,1077.4)--(764.9,1077.4);
\draw[hdgrid] (137,1102)--(764.9,1102);
\draw[hdgrid] (137,1126.6)--(764.9,1126.6);
\draw[hdoutline] (281.9,757.6)--(330.2,806.8)--(281.9,856)--(233.6,806.8)--cycle;
\node[font=\normalsize,inner sep=0pt] at (293.975,816.64) {A};
\draw[hdoutline] (547.55,1028.2)--(595.85,1077.4)--(644.15,1028.2)--(644.15,1102)--(547.55,1102)--cycle;
\node[font=\normalsize,inner sep=0pt] at (607.925,1089.7) {B};
\draw[hdoutline] (185.3,954.4)--(185.3,905.2)--(281.9,905.2)--(233.6,954.4)--cycle;
\node[font=\normalsize,inner sep=0pt] at (223.94,942.1) {C};
\draw[hdoutline] (450.95,806.8)--(475.1,831.4)--(475.1,880.6)--(450.95,905.2)--(426.8,880.6)--(426.8,831.4)--cycle;
\node[font=\normalsize,inner sep=0pt] at (463.025,868.3) {D};
\fill[black] (185.3,954.4) circle (1.7);
\node[font=\normalsize,inner sep=0pt] at (171,970) {$O$};
\fill[black] (450.95,806.8) circle (1.7);
\node[font=\normalsize,inner sep=0pt] at (460,794) {$O\prime$};
\end{scope}

\begin{scope}[color=answerred,draw=answerred,text=answerred]\draw[answerred,line width=1pt] (274.07,265.40) -- (290.37,232.47) -- (355.57,232.47) -- (339.27,265.40) -- cycle;
\draw[answerred,line width=1pt] (412.67,587.33) -- (444.67,620.00) -- (412.67,653.33) -- (380.00,620.00) -- cycle;
\draw[answerred,line width=1pt] (219.33,604.00) -- (251.33,636.67) -- (283.33,604.00) -- (283.33,653.33) -- (219.33,653.33) -- cycle;
\draw[answerred,line width=1pt] (122.67,636.67) -- (154.67,636.67) -- (154.67,702.00) -- (122.67,669.33) -- cycle;
\draw[answerred,line width=1pt] (299.33,538.00) -- (315.33,522.00) -- (347.33,522.00) -- (363.33,538.00) -- (347.33,554.67) -- (315.33,554.67) -- cycle;
\ans{405.33}{612.67}{26.00}{11}{$A\prime$}
\ans{242.00}{636.67}{26.00}{11}{$B\prime$}
\ans{133.33}{670.67}{26.00}{11}{$C\prime$}
\ans{328.67}{531.33}{26.00}{11}{$D\prime$}
\end{scope}
\end{scope}\pgfresetboundingbox\path[use as bounding box] (86.0,483.0) rectangle (515.0,756.0);
\end{tikzpicture}}\par

\PlainItem{① 将图A向右平移14格后，再向下平移5格。}
\par

\PlainItem{② 将图B向左平移9格后，再向上平移5格。}
\PlainItem{③ 将图C绕$O$点顺时针旋转$90^\circ$。}
\PlainItem{④ 将图D绕$O\prime$点逆时针旋转$90^\circ$。}




\QHeading{旋转90度的方格法}

\Answer{以旋转中心为起点记录顶点的横向、纵向格数。顺时针旋转时，原来“向右a格、向上b格”的点变成“向右b格、向下a格”；逆时针则变成“向左b格、向上a格”。这里a、b表示所数格数，不要求学生建立坐标方程。其他方向可以先用手势转向核对；原点到顶点的距离必须不变。}

\QHeading{平移的另一种检查办法}

\Answer{可以分别对照对应边的长度和方向，再数任一对对应顶点之间的横、竖格数；也可以逐顶点平移再顺次连线。图中每个顶点必须移动相同方向、相同距离，不能只移动其中一个点后改变图形。}
